\subsection{Normal subgroups under two-dimensional homology extensions}

\begin{claim}{A.23}{Proved}{Independent relative exponent rows preserve every coefficient quotient}
Let $H$ be a group, let $S$ and $R$ be sets, and let
$w_r\in H*F(S)$ for $r\in R$. Suppose that the vectors of exponent
sums of the letters in $S$ in the words $w_r$ are linearly
independent over $\mathbb Q$. Set
\[
 E=(H*F(S))/\langle\!\langle w_r:r\in R\rangle\!\rangle.
\]
For every group homomorphism $u:H\to K$, the canonical map
\[
 K\longrightarrow K*_H E
\]
is injective. In particular, $H\to E$ is injective, and every
normal subgroup $N\triangleleft H$ satisfies
\[
 \langle\!\langle N\rangle\!\rangle_E\cap H=N.
 \tag{A.23.1}
\]
Here $K*_H E$ is the group pushout of $u$ and the canonical
map $H\to E$; no injectivity of $u$ is assumed.

This applies to the map on fundamental groups of every inclusion
$Y\subset X$ of connected CW complexes of dimension at most two
that is an ordinary integral homology isomorphism. If $X$ is
aspherical, then $Y$ is aspherical too. If additionally
$\pi_1(X)$ is simple and $\pi_1(Y)\ne1$, then $\pi_1(Y)$ is simple.
Consequently such a two-dimensional homology extension cannot
place an old element $h\in H$ in the normal closure of an old
$g\in H$ unless $h$ already belonged to that normal closure in $H$.
\end{claim}
\begin{proof}
For each $r$, apply $u$ to the coefficient letters of $w_r$ and
leave the variable letters fixed; write the resulting word as
$u_*w_r\in K*F(S)$. The universal property of a presented group
gives an explicit identification
\[
 K*_H E\ \cong\
 (K*F(S))/\langle\!\langle u_*w_r:r\in R\rangle\!\rangle:
 \tag{A.23.2}
\]
both sides universally receive $K$ and one element for each
$s\in S$, subject precisely to these displayed relations.
The relative exponent rows have not changed. The nonsingular
coefficient-injectivity theorem
\cite[Theorem~1 and Remark~5.1]{D_OpenAI2026Nonsingular}, checked
in the source audit and used explicitly in A.19, therefore makes
the canonical map from $K$ injective. The infinite-set assertion
can also be reduced directly to its finite theorem: a coefficient
killed in (A.23.2) has a finite normal-closure expression, involving
finitely many relators and finitely many variable letters. Those
relators still have independent exponent rows, and their entire
finite supports can be retained. The finite theorem says that the
coefficient was already trivial in $K$.

Taking $u=1_H$ proves the injection $H\to E$. Taking
$K=H/N$ and $u$ the quotient gives
\[
 K*_H E\ \cong\ E/\langle\!\langle N\rangle\!\rangle_E.
\]
Since $H/N$ maps injectively to this group, the kernel of the
composite from $H$ is exactly $N$. It is also
$\langle\!\langle N\rangle\!\rangle_E\cap H$, proving
(A.23.1).

For the topological application, choose a maximal tree $T_Y$ of
$Y^{(1)}$ and extend it to a maximal tree $T_X$ of $X^{(1)}$.
The fact that $T_Y$ already connects every vertex of $Y$ gives
$T_X\cap Y=T_Y$. Collapse these trees. The resulting pair
$\overline Y\subset\overline X$ has one common vertex and is
homotopy equivalent to the original pair on each individual
space. Its relative cellular complex consists of free groups in
degrees two and one, and its homology is zero. Thus its boundary
\[
 C_2(\overline X,\overline Y;\mathbb Z)
 \xrightarrow{\ \cong\ }
 C_1(\overline X,\overline Y;\mathbb Z)
\]
is an isomorphism. Present $\pi_1(X)$ over $H=\pi_1(Y)$ by its
relative one-cells and attaching words of the relative two-cells.
Their exponent rows are precisely the columns of this cellular
boundary after choosing bases, hence are independent. This proves
the asserted injection on fundamental groups and its universal
pushout and normal-subgroup conclusions.

When $X$ is aspherical, the proved injection identifies the
universal cover of $Y$ with a component of the preimage of $Y$
in the contractible universal cover of $X$. It is a
simply connected two-dimensional subcomplex. Its two-cycles
inject into the cellular two-cycles upstairs in $X$, which are
zero because that cover is contractible and has no three-cells.
It therefore has vanishing reduced homology in every degree;
Hurewicz and Whitehead make it contractible. This proves
asphericity of $Y$.

Finally let $E=\pi_1(X)$ be simple and $1\ne H=\pi_1(Y)$.
For any nontrivial normal subgroup $N$ of $H$, its normal closure
in $E$ is a nontrivial normal subgroup of $E$, hence equals $E$.
Equation (A.23.1) now says $N=H$, proving that $H$ is simple.
Taking $N=\langle\!\langle g\rangle\!\rangle_H$ proves the
last normal-generation assertion, including the case $g=1$.
\end{proof}
